\bookchapter{第五十八章：速查手冊同資源}

\chapterintro{呢一章係一份可以獨立使用嘅速查手冊。當你記唔清某個寫法或者需要快速決策嗰陣，直接揭到嚟呢度。建議當佢係「口袋卡片」：常用代碼速查、常見問題、學習資源，一頁都唔漏。}

\booksection{一、常用代碼速查卡}

\begin{center}
\begin{tabularx}{\linewidth}{|X|X|}
\hline
\textbf{要做乜嘢} & \textbf{關鍵代碼} \tabularnewline
\hline
記住一個可變狀態 & var x by remember \{ mutableStateOf(初值) \} \tabularnewline
轉屏都唔會冇咗狀態 & rememberSaveable \{ mutableStateOf(初值) \} \tabularnewline
直排/橫排佈局 & \texttt{Column \{ \}}/\texttt{Row \{ \}} \tabularnewline
懶加載列表 & LazyColumn \{ items(列表) \{ ... \} \} \tabularnewline
頁面跳轉 & navController.navigate("路由") \tabularnewline
喺協程入面做要等嘅嘢 & \texttt{viewModelScope.\allowbreak launch \{ .\allowbreak .\allowbreak .\allowbreak  \}} \tabularnewline
觀察數據流 & \texttt{val list by flow.\allowbreak collectAsState()} \tabularnewline
儲存小設定 & context.dataStore.edit \{ it[KEY] = 值 \} \tabularnewline
存結構化數據 & Room 三件套：@Entity / @Dao / @Database \tabularnewline
聯網攞數據 & Retrofit 接口 + suspend 函數 \tabularnewline
發通知 & NotificationChannel + NotificationCompat.Builder \tabularnewline
打日誌 & \texttt{Log.\allowbreak d("TAG", "message")} \tabularnewline
\hline
\end{tabularx}
\end{center}

\begin{center}\small\color{gray}安卓開發速查卡（2026 年 8 月版）\end{center}

\booksection{二、常見問題 FAQ}

Q1：我完全唔識編程，學安卓開發要幾耐？

跟本書嘅路線行：第 1-2 周裝環境、學 Kotlin 基礎；第 3-4 周學 Compose 做到界面出嚟；第 5-6 周完成記賬本項目；第 7-8 周做完天氣項目。兩個月，由零到識得獨立做小 App。前提係每日投入 1-2 個鐘，而且要動手敲代碼，唔係齋睇書。

Q2：Java 仲使唔使學？

唔使。新項目、新教程、新崗位基本上都以 Kotlin 為準。只有「維護十年老項目」嗰陣先會遇到大量 Java。學完 Kotlin 之後如果工作需要，再補 Java 會好快——兩者底層係同一個虛擬機生態。

Q3：模擬器定係真機？

開發除錯用模擬器方便（有快照、多機型），但驗收一定要過真機——觸感、性能、真網絡、感應器，模擬器代替唔到。條件許可嘅話，買一部中端真機做測試機，係最劃算嘅投資。

Q4：安卓開發容唔容易搵到工？

2026 年嘅行情：齋「識寫界面」嘅初級崗位競爭激烈，但「識 Kotlin + 懂架構 + 有上線作品」嘅開發者依然搶手。學習期嘅你，最好嘅履歷就係商店入面嗰個自己做嘅 App。

Q5：遇到報錯卡住咗點算？

黃金三步：(1) 將報錯第一行原原本本攞去搜 StackOverflow/搜尋引擎；(2) 查官方文檔對應頁面；(3) 仲係唔得，就將「期望行為 + 實際現象 + 報錯全文 + 相關代碼」打包發畀 AI 或者社區發問。絕大多數問題，第一步就搞掂。

\booksection{三、推薦學習資源}

\begin{center}
\begin{tabularx}{\linewidth}{|X|X|X|}
\hline
\textbf{類型} & \textbf{資源} & \textbf{推薦理由} \tabularnewline
\hline
官方 & developer.android.com & 安卓官方文檔，Compose/Kotlin 教程都喺入面，第一權威 \tabularnewline
官方 & Kotlin 官網 (kotlinlang.org) & 語言官方文檔 + 網上演練場 \tabularnewline
官方 & Google Codelabs & 手把手嘅分步實驗課，新手友好度滿分 \tabularnewline
社區 & Stack Overflow (android 標籤) & 報錯求救第一站 \tabularnewline
社區 & 掘金 / CSDN 安卓專區 & 中文實戰文章多，注意分辨新定舊（先睇日期） \tabularnewline
影片 & Google Developers 官方頻道 & 每年嘅開發者大會新技術講解 \tabularnewline
配套 & 第七篇《AI編程實戰篇》（第37-47章） & 用 AI 輔助安卓開發嘅完整心法 \tabularnewline
\hline
\end{tabularx}
\end{center}

\begin{center}\small\color{gray}安卓開發學習資源清單（2026 年 8 月）\end{center}

\booksection{四、Gradle 常用命令速查}

平時㩒掣就夠，但命令行有時更高效（尤其係配合 AI 工具或者 CI 流水線嗰陣）：

\begin{center}
\begin{tabularx}{\linewidth}{|X|X|}
\hline
\textbf{命令} & \textbf{作用} \tabularnewline
\hline
\texttt{.\allowbreak /\allowbreak gradlew assembleDebug} & 打一個除錯版 APK（Windows 記得將 ./ 攞走） \tabularnewline
\texttt{.\allowbreak /\allowbreak gradlew assembleRelease} & 打發布版（要配好簽名） \tabularnewline
\texttt{.\allowbreak /\allowbreak gradlew test} & 行晒全部單元測試 \tabularnewline
\texttt{.\allowbreak /\allowbreak gradlew lint} & 靜態檢查：搵出潛在 bug 同唔規範嘅寫法 \tabularnewline
\texttt{.\allowbreak /\allowbreak gradlew clean} & 清空構建產物（「重啟大法」嘅編譯版） \tabularnewline
\texttt{.\allowbreak /\allowbreak gradlew :app: dependencies} & 睇依賴樹（排查版本衝突神器） \tabularnewline
\hline
\end{tabularx}
\end{center}

\begin{center}\small\color{gray}Gradle 常用命令速查\end{center}

\begin{tcbinfo}{小貼士：構建出問題先 clean}
「尋日仲行得今日又唔得」嘅靈異事件，十件有六七件係構建緩存錯亂。./gradlew clean 之後重新構建，可以醫好一大半玄學問題。仲唔得就重啟 Android Studio——呢招請配合 File Invalidate Caches 一齊服用。

\end{tcbinfo}

\booksection{五、結語：由第一個 App 到第 N 個 App}

\booksection{七、結語：由第一個 App 到第 N 個 App}

返轉頭睇，你由「唔知 Activity 係乜嘢」，到裝好環境、識咗 Kotlin、用 Compose 畫出界面、畀 App 接上數據庫同網絡，再到測試、打包、上架——呢條路上你已經行完咗由零到一最難嗰一段。

技術嘅世界變得好快：框架會升級，工具會換代，AI 會接手越嚟越多「敲鍵盤」嘅嘢。但有一樣嘢唔會過時：你將一個朦朧嘅想法，拆解成清晰嘅步驟，再一步一步實現佢嘅能力。工具係放大器，思路係地基。放大器可以換，地基唔可以鬆動。

你唔再只係"識玩手機嘅人"。你可以將腦入面嘅諗法，變成喺三億部安卓裝置上面真實運行嘅 App。由今日開始，用鍵盤去創造——呢個就係安卓開發者嘅浪漫。

\booksection{參考文獻}

\begin{enumerate}
  \item Android Developers. (2026). Android Developer Guides and Jetpack Compose Documentation. developer.android.com.
  \item Kotlin Foundation. (2026). Kotlin Language Documentation. kotlinlang.org.
  \item Android Open Source Project. (2026). Android Platform Version History and API Levels.
  \item Google Play. (2026). Play Console Help: Publishing and App Signing.
  \item AndroidX Release Notes. (2026). Room / DataStore / Navigation / WorkManager releases.
  \item Square. (2026). Retrofit Documentation. github.com/square/retrofit.
\end{enumerate}
